\begin{lemma}[Bad-witness characterization]
Let $r$ be a quasi-order on $Q$.  The domination quasi-order on
$\mathcal P(Q)$ fails to be well-quasi-ordered if and only if there is a
bad sequence of sequences $f:[\omega]^2\to Q$, namely
\[
\forall m<n<l\qquad f(m,n)\nleq f(n,l).
\]
\end{lemma}
\begin{proof}
Assume first that the relation from \noderef{SetDomination},
$\mathsf{SetDomination}(r)$, is not well-quasi-ordered.
By the sequence characterization of well-quasi-ordering, there is a
sequence $(P_n)_{n\in\omega}$ of subsets of $Q$ for which no $m<n$ has
$P_m\preceq_r P_n$.  For each $m<n$, failure of domination supplies
$q_m^n\in P_m$ such that no member of $P_n$ is $r$-above $q_m^n$.
Define $f(m,n)=q_m^n$.  If $m<n<l$ and
$r(f(m,n),f(n,l))$ held, then $f(n,l)\in P_n$ would contradict the
choice of $q_m^n$.  Hence $f$ is bad in the sense of
\noderef{BadPairSequence}.

Conversely, let $f$ be bad and put
\[
P_m=\{f(m,n)\mid m<n\}.
\]
If $m<n$, then $f(m,n)\in P_m$.  It cannot be dominated by $P_n$:
every element of $P_n$ has the form $f(n,l)$ with $n<l$, and badness
for $m<n<l$ says that
$\neg r(f(m,n),f(n,l))$.  Thus no $P_m$ is dominated by a later
$P_n$, so this sequence witnesses failure of
$\mathsf{WellQuasiOrdered}(\mathsf{SetDomination}(r))$.
The two implications prove the equivalence.
\end{proof}
